\section{Evaluation}
\label{sec:evaluation}
We evaluate inference coverage, optimizer integration,
specification size, and compilation and execution costs. Coverage results
use one frozen optimized compiler; size, cost, and XLA integration
experiments retain their separate snapshots. TEPL proofs
use algebraic semantics rather than strict floating-point laws.
Appendix~\ref{app:current-evaluation} records protocols and provenance.

\sidefigure[0.4]{figures/XLA_rewritie_cases.pdf}
{XLA case coverage and attribute synthesis result. 31 out of 38 cases with RHS axis field are inferred}
{fig:xla-rewrite-cases}
{The inventory separates modeled and unmodeled cases and reports
acceptance for axis-inference and verification-only cases.}

\subsection{Inference Coverage}
We evaluate attribute synthesis on 209 concrete XLA rewrite
cases~\cite{xla-evaluation-source}. A case is \emph{modeled} when
TEPL's verifier supports the semantics of all operations it requires.
Of the 209 cases, 123 are modeled and 86 are unmodeled.
Unmodeled does not mean that a rewrite is inapplicable.
These rules can still be expressed and applied with explicitly supplied
attributes, but TEPL cannot verify them or infer their axis attributes
with its current semantic support.

Of the modeled cases, 38 are \emph{axis-inference cases}, which require
RHS axis attributes to be inferred. The other 85 are
\emph{verification-only cases}: their replacements use supplied attributes,
attribute-free operations such as elementwise addition, or directly
return a matched value, as in $\operatorname{add}(X,0)\to X$.
Figure~\ref{fig:xla-rewrite-cases} summarizes the results.

A case is \emph{accepted} when semantic proof, compilation, and
validation succeed with the original guards and non-axis fields
unchanged. TEPL accepts 98 of 123 modeled cases (79.7\%),
including 67 of 85 verification-only cases.
Attribute synthesis succeeds on 31 of 38 axis-inference cases
(81.6\%), inferring all requested RHS axis attributes while preserving
the written rewrite conditions.

\subsection{Integration into an Existing XLA Pass}
\label{sec:evaluation-xla-integration}
We integrate TEPL into XLA's existing \code{AlgebraicSimplifier}
pass~\cite{xla-evaluation-source}. From compiled TEPL rules, we generate
adapter code and C++ pattern-matching and attribute-transformation logic.
This reuses TEPL's code-generation infrastructure while allowing
control over how the generated code is embedded in XLA.

All rewrite rules in the integrated pass are expressed in TEPL,
including unmodeled rules. RHS attributes are inferred where supported;
attribute transformations for unmodeled rules are written explicitly
in TEPL.

We evaluate four representative workloads. For each workload, we count
the distinct XLA rewrite rules applied, their total applications,
and how many applied rules are modeled
(Table~\ref{tab:algebraic-simplifier-workloads}).
The integrated rewrite logic passes all 28 related XLA test cases.
For all four workloads, the final compiled graphs have the same
structure as those produced by the unmodified XLA pass, and the
output values agree.

\begin{table}[!htbp]
\caption{Rewrite applications in four representative workloads.}
\label{tab:algebraic-simplifier-workloads}
\small
\begin{tabular*}{\linewidth}{@{\extracolsep{\fill}}lrrrc@{}}\toprule
Workload & Applied rules & Total firings & Modeled rules & Final graph match\\\midrule
DLRM training & 32 & 201 & 20/32 & Yes\\
GPT-2 Small & 14 & 1,693 & 6/14 & Yes\\
ViT & 17 & 1,850 & 9/17 & Yes\\
ResNet-50 & 10 & 1,027 & 5/10 & Yes\\\bottomrule
\end{tabular*}
\end{table}

\subsection{Shorter Rewrite Rules}
\label{sec:evaluation-size}
We compare code length across three implementations of the 31 rules
for which axis inference succeeds: XLA, TEPL with explicit attribute
derivations, and TEPL with inferred axis attributes.
We measure source size in lexical tokens, excluding comments,
whitespace, imports, and shared declarations.
Table~\ref{tab:specification-size} reports the results.

\begin{table}[!htbp]
\caption{Source footprint of all 31 inference-successful rules.
Counts exclude comments, whitespace, imports, and shared declarations.}
\label{tab:specification-size}
\small
\definecolor{specificationBlue}{HTML}{3F608A}
\definecolor{specificationGreen}{HTML}{396F59}
\setlength{\tabcolsep}{4pt}
\begin{tabularx}{\linewidth}{@{}Yrrr@{}}\toprule
Section & XLA & Explicit TEPL\textsuperscript{a} & Axis inferred TEPL\textsuperscript{b}\\\midrule
Matching (LHS) & 1,291 & \textcolor{specificationBlue}{(-70.2\%)}385 & \textcolor{specificationGreen}{(-4.2\%)}369\\
Legality checks (\code{where} block) & 1,226 & \textcolor{specificationBlue}{(-25.0\%)}920 & \textcolor{specificationGreen}{(-0.0\%)}920\\
Attribute reasoning (\code{derive} block) & 4,915 & \textcolor{specificationBlue}{(-75.1\%)}1,225 & \textcolor{specificationGreen}{(-62.3\%)}462\\
Replacement construction (RHS) & 1,668 & \textcolor{specificationBlue}{(-80.8\%)}320 & \textcolor{specificationGreen}{(-16.2\%)}268\\
Other syntax & 922 & \textcolor{specificationBlue}{(-61.5\%)}355 & \textcolor{specificationGreen}{(-0.0\%)}355\\\midrule
Total & 10,022 & \textcolor{specificationBlue}{(-68.0\%)}3,205 & \textcolor{specificationGreen}{(-25.9\%)}2,374\\\bottomrule
\end{tabularx}
\par\smallskip
{\footnotesize\raggedright
\textsuperscript{a}\textcolor{specificationBlue}{Blue percentages} show the decrease from XLA to Explicit TEPL.
\textsuperscript{b}\textcolor{specificationGreen}{Green percentages} show the decrease from Explicit to Inferred TEPL.\par}
\end{table}

Explicit TEPL uses 68.0\% fewer tokens than XLA.
Axis inference reduces the TEPL token count by a further 25.9\%.
Most of this additional reduction comes from attribute derivations,
whose token count decreases by 62.3\%.
The comparison with XLA is not direct, since its implementation
also includes logic that TEPL does not support.
However, the reduction from explicit to inferred TEPL shows how
automatic inference can reduce the code needed for manual attribute
reasoning.

\subsection{Inference Scalability}
\label{sec:evaluation-cost}
Mean compilation time grows approximately linearly and stays below 100\,ms
for both 1--64 dots per side and 1--64 additive terms. One compiled rule
handles varying tensor ranks without recompilation. Mean runtime metadata
program execution time grows with rank and stays below 1\,ms at rank 128.
Appendix~\ref{app:evaluation-scalability} gives the experimental setup and
detailed results.
